\documentclass[11pt]{article}
\usepackage{fontspec}
\usepackage[french]{babel}
\usepackage{amsmath}
\usepackage{amssymb}
\usepackage{amsthm}
\usepackage{mathrsfs}
\usepackage{tikz-cd}
\usetikzlibrary{babel}
\usepackage{enumitem}

\newcommand{\printedpage}[1]{\marginpar{\footnotesize p.~#1}}
\newcommand{\msp}[1]{\marginpar{\footnotesize [#1]}}
% underlined category names, as printed (e.g. \uc{Fibhot}\,S)
\newcommand{\uc}[1]{\underline{\mathrm{#1}}}
\newcommand{\HOT}{\mathrm{HOT}}
\newcommand{\Cat}{\mathrm{Cat}}
\newcommand{\defeq}{\overset{\text{déf}}{=}}

\begin{document}

\printedpage{29}
% head: A. Grothendieck : Dérivateurs, Ch. VI, Édition des Universités de Montpellier II et Paris VII
% foot: Version 9 novembre 2001

\begin{enumerate}[leftmargin=*]
\item[d)] Pour que $f : X \longrightarrow Y$ soit dans $W$, il faut et il suffit que $\pi_0(f)$ soit bijectif, et que les $f_i : X_i \longrightarrow Y_i$ induits sur les composantes connexes soient dans $W$.
\end{enumerate}

Ces axiomes sont vérifiés par tous les $W$ qui je connais, et qu' on aurait envie de regarder comme localiseurs. Par contre, L5, L6, L6 bis (p. 43, 44) sont des axiomes de nature beaucoup plus spéciale. J'ai envie de dire qu'ils distinguent les localiseurs de nature proprement \textit{`homotopique'} de ceux qui sont de nature \textit{`homologique'} (ou `cohomologique'). Il faudrait que je m'assure qu'ils sont satisfaits pour les $W_i$ ($i$-équivalence d'homotopie) (\footnote[32]{C'est \textit{faux}, cf. annotation ci-dessus.}). % [sic] qui je connais -> que je connais, relative pronoun object requires "que"

Le fait que L4 bis, i.e. c) ci-dessus, implique l'axiome dual, résulte du fait que L4 bis implique ceci

\textbf{Proposition 2.} \textit{Soit $f : X \longrightarrow Y$ dans} $\Cat$\textit{, alors $f \in W$ si et seulement si $f^o : X^o \longrightarrow Y^o$ est dans $W$.}

Se voit par l'argument de I, page 75 (cor. 2) - n'utilise que L4, pas même L1, L3.

\msp{page 46}

Je vais essayer de résumer ce que je crois avoir compris :

\begin{enumerate}[leftmargin=*]
\item[1)] La $W$-équivalence `par fibres' n'a d'intérêt (pour construire des catégories localisées) que dans le cas où on sait qu'elle implique déjà, soit la $W$-équivalence locale, soit la $W$-équivalence colocale (soit même la $W$-équivalence universelle) sur $S$, donc quand on suppose $X$, $Y$ pour le moins tous deux soit $W$-lisses, soit $W$-propres, soit $W$-fibrés sur $S$.

\item[2)] On pourrait appeler la $W$-équivalence locale sur $S$ la \textit{`$W$-équivalence fine à gauche'} sur $S$, et la $W$-équivalence colocale sur $S$ la \textit{`$W$-équivalence fine à droite'} sur $S$. Elles coïncident dans $\uc{Fibhot}\,S$, a fortiori dans $\uc{Parf}\,S$ et dans les quatre catégories (cf. diagramme p. 8) $\uc{Fib}_0S$, $\uc{Fibsc}_0S$, $\uc{Cofib}_0S$, $\uc{Cofibsc}_0S$ (\footnote[33]{On pourrait les appeler les six catégories \textit{internes} du diagramme.}). Par contre, elles ne coïncident pas sur les catégories \textit{bordantes} du diagramme, notamment sur $\uc{Liss}\,S$, $\uc{Prop}\,S$, $\uc{Fib}\,S$, $\uc{Cofib}\,S$, $\uc{Fibsc}\,S$, $\uc{Cofibsc}\,S$. À cet égard, c'est la \textit{$W$-équivalence fine} \textbf{droite} (en termes des $s\backslash X$) qui est intéressante sur le côté (hélas) gauche du diagramme, i.e. pour $\uc{Liss}\,S$, $\uc{Fib}\,S$, $\uc{Fibsc}\,S$, alors que c'est la \textit{$W$-équivalence fine} \textbf{gauche} (en termes des $X/s$) qui est pertinente du côté droit, pour $\uc{Prop}\,S$, $\uc{Cofib}\,S$, $\uc{Cofibsc}\,S$. On n'aura jamais à regarder simultanément les deux équivalences fines sur une même catégorie du diagramme, à la seule exception de $\Cat/S$ tout entier.
\end{enumerate}

\msp{page 47}

\begin{enumerate}[leftmargin=*]
\item[3)] La $W$-équivalence grossière, elle, est `bonne partout' - mais ça dépend pour qui. Pour les 6 catégories qui sont au dessus de $\uc{Fibhot}\,S$, i.e.
\[
\uc{Fibhot}\,S,\quad \uc{Parf}\,S,\quad \uc{Fib}_0S,\quad \uc{Fibsc}_0S,\quad \uc{Cofib}_0S,\quad \uc{Cofibsc}_0S,
\]
\printedpage{30}
% head: A. Grothendieck : Dérivateurs, Ch. VI, Édition des Universités de Montpellier II et Paris VII
% foot: Version 9 novembre 2001
elle coïncide avec toutes les autres équivalences (moyennnant l'axiome L5), pour les trois catégories bordantes supérieures gauches % [sic] moyennnant -> moyennant, triple n
\[
\uc{Liss}\,S,\quad \uc{Fib}\,S,\quad \uc{Fibsc}\,S
\]
du diagramme, elle coïncide avec la $W$-équivalence fine droite, et dualement sur
\[
\uc{Prop}\,S,\quad \uc{Cofib}\,S,\quad \uc{Cofibsc}\,S
\]
elle coïncide avec la $W$-équivalence fine gauche. Quant à $\Cat/S$ tout entier, c'est la seule des catégories du diagramme où il y ait \textit{trois} notions intéressantes différentes de $W$-équivalence, à savoir : la grossière, et les deux fines (à gauche, à droite). Sur les autres six catégories bordantes du diagramme, il y en a en principe deux : la grossière, et une spécifiée des deux fines (défendu de regarder l'autre - à ses risques et périls \ldots) : la droite du côté gauche, la gauche du côté droit. (NB Je devrais refaire le diagramme, décidément, en mettant les fibrations à droite, pas à gauche.) Quant aux six catégories internes au diagramme, dans elles il n'y a plus qu'une seule notion de $W$-équivalence (puisque sur elles les quatre notions coïncident).
\end{enumerate}

\msp{page 48 (vide)}

\msp{page 49} (\textbf{NB} il n'y a pas de page 48)

Je résume maintenant ce que je crois avoir compris sur les catégories localisées correspondantes. Il y en a, à équivalence près, \textit{trois} différentes seulement en tout et pour tout, quand on prend les localiseurs `admissibles' dans les catégories, comme je viens de le dire (c'étant ça en somme l'intérêt objectif des prescriptions!) De façon précise :

\begin{enumerate}[leftmargin=*]
\item[$1^o$)] La \textit{$W$-équivalence grossière} donne `la même' catégorie localisée pour les 13 catégories du diagramme, i.e. toutes les flèches de transition du diagramme induisent des équivalences pour cette $W$-équivalence grossière. Le localisé sera noté \fbox{$\HOT^{\mathrm{lc}}_W(S)$}, où l'exposant lc signifie: ($W$-types d'homotopie relatif sur $S$) \textit{`localement constants.'} (Il ne me reste qu'un petit doute pour $\uc{Fibsc}\,S$ et $\uc{Cofibsc(S)}$ (\footnote[34]{Ce doute n'est pas fondé, tout marche parfaitement!}), si l'équivalence grossière y est bel et bien raisonnable. On verra en faisant la démonstration.)

\item[$2^o$)] La \textit{$W$-équivalence fine à droite} donne le même résultat pour toutes les catégories du diagramme où elle est considérée comme admise, \textit{et où elle ne coïncide pas avec la $W$-équivalence grossière} déjà considérée, donc sur les quatre catégories bordantes gauches
\[
\Cat/S,\quad \uc{Liss}\,S,\quad \uc{Fib}\,S,\quad \uc{Fibsc}\,S
\]
(\footnote[35]{\textbf{NB} Sur les trois dernières cette équivalence est la $W$-équivalence par fibres.}). Le localisé est noté \fbox{$\HOT_W(S)$}.
\end{enumerate}

\msp{page 50}

\printedpage{31}
% head: A. Grothendieck : Dérivateurs, Ch. VI, Édition des Universités de Montpellier II et Paris VII
% foot: Version 9 novembre 2001

\begin{enumerate}[leftmargin=*]
\item[$2'^o$)] Dualement, la \textit{$W$-équivalence fine gauche} donne le même résultat sur toutes les catégories du diagramme où elle est admise, et où elle ne coïncide pas avec la grossière, i.e. sur les 4 catégories
\[
\Cat/S,\quad \uc{Prop}\,S,\quad \uc{Cofib}\,S,\quad \uc{Cofibsc}\,S.
\]
Il y a lieu de noter la catégorie des fractions commune par le sigle
\[
\HOT_W(S^o)^o \qquad \text{(ou aussi $\HOT^o_W(S)$)},
\]
puisqu'elle est isomorphe à celle notée par ce sigle.
\end{enumerate}

Si on veut des notations sans ambiguïté, désignant des catégories bien définies, et non pas seulement `à équivalence près', je choisirai
\[
\HOT_W(S) \defeq \uc{Fibsc}(S)(W^{\mathrm{d}}_S)^{-1}
\]
d'où
\[
\HOT_W(S^o) \qquad \underbrace{\simeq}_{\text{\footnotesize isomorphisme de catégories}} \qquad \uc{Cofibsc}(S)(W^{\mathrm{g}}_S)^{-1},
\]
où $W^{\mathrm{d}}_S$, $W^{\mathrm{g}}_S$ désignent les ensembles de flèches correspondants à la $W$-équivalence fine droite resp. gauche. Pour $\HOT^{\mathrm{lc}}_W(S)$, je préfère une définition aussi symétrique que possible en $S$, $S^o$, et de plus une catégorie aussi petite que possible, ce qui me conduit à choisir
\[
\HOT^{\mathrm{lc}}_W(S) \defeq (\uc{Parf}\,S)W^{-1}_S,
\]
où on a noté par $W_S$ les $W$-équivalences dans $\uc{Parf}\,S$. (Inutile cette fois de préciser lesquelles - toutes les quatre \texttt{[notions]} coïncident dans $\uc{Parf}\,S$ !)

\msp{page 51}

Notons l'isomorphisme
\[
\uc{Parf}(S^o) \simeq \uc{Parf}\,S
\]
(provenant du fait que $f : X \longrightarrow Y$ dans $\Cat$ est une fibration parfaite si et seulement si $f^o$ l'est). Cet isomorphisme $X \longmapsto X^o$ est compatible avec les localiseurs (car $f \in W \Longleftrightarrow f^o \in W$), d'où
\[
\HOT^{\mathrm{lc}}_W(S^o) \simeq \HOT^{\mathrm{lc}}_W(S)
\]
(isomorphisme de catégories). On aura deux foncteurs canoniques entre les trois catégories de $W$-types d'homotopie relatifs sur $S$
\[
\begin{tikzcd}[column sep=small]
(\uc{Fibsc}_0S)W_S^{-1} \arrow[rr, "\approx"] \arrow[dr, dotted, "i'"'] & & \HOT^{\mathrm{lc}}_W(S) \arrow[dl, "i"] \arrow[dr, "j"'] & & (\uc{Cofibsc}_0S)W_S^{-1} \arrow[ll, "\approx"'] \arrow[dl, dotted, "j'"] \\
& \HOT_W(S) & & \HOT^o_W(S) \defeq \HOT_W(S^o). &
\end{tikzcd}
\]
\printedpage{32}
% head: A. Grothendieck : Dérivateurs, Ch. VI, Édition des Universités de Montpellier II et Paris VII
% foot: Version 9 novembre 2001
\texttt{[plutôt} $\defeq \HOT_W(S^o)^o$\texttt{]} Ces flèches $i$, $j$ sont suggerées par les flèches en pointillés $i'$, $j'$, provenant des inclusions % [sic] suggerées -> suggérées, missing accent
\[
\uc{Fibsc}_0S \lhook\joinrel\longrightarrow \uc{Fibsc}\,S,\quad \uc{Cofibsc}_0S \lhook\joinrel\longrightarrow \uc{Cofibsc}\,S,
\]
qui définiraient $i$, $j$ à l'aide du choix de foncteurs quasi-inverses pour les flèches horizontales. Mais $i$, $j$ ne seraient définis qu'à isomorphisme (unique) près. Il sera plus joli de choisir les quasi-inverses de façon canonique, à l'aide des flèches
\[
\uc{Fibsc}_0S \xleftarrow{\Phi_0} \uc{Parf}\,S \xrightarrow{\Psi_0} \uc{Cofibsc}_0S
\]
(où $\Phi_0 = \Phi|\uc{Parf}\,S$, $\Psi_0 = \Psi|\uc{Parf}\,S$) du diagramme p. 8. La question principale
\msp{page 52}
qui reste non élucidée, c'est si ces foncteurs $i$, $j$ sont bien \textit{pleinement fidèles,} comme je m'y attends certes.

Je veux maintenant prouver les affirmations précédentes sur les localisés.

\textcircled{1} Le localisé pour la $W$-équivalence grossière. Il suffira, à un $\varepsilon$ près, d'établir les conditions du lemme p. 11, dans les trois cas suivants.
\begin{tikzcd}[baseline=(\tikzcdmatrixname-1-1.base)]
\uc{Parf}\,S \arrow[r, "r"] & \Cat/S \arrow[l, dotted, bend left, "C"]
\end{tikzcd}
(ça donnera le cas de $\uc{Fibhot}\,S$, $\uc{Liss}\,S$ et $\uc{Prop}\,S$ sur $\Cat/S$) et
\begin{tikzcd}[baseline=(\tikzcdmatrixname-1-1.base)]
\uc{Fibsc}_0S \arrow[r, "p"] & \Cat/S \arrow[l, dotted, bend left]
\end{tikzcd},
\begin{tikzcd}[baseline=(\tikzcdmatrixname-1-1.base)]
\uc{Cofibsc}_0S \arrow[r, "q"] & \Cat/S \arrow[l, dotted, bend left]
\end{tikzcd},
qui donneront respectivement les cas de
\[
\uc{Fib}_0S,\ \uc{Fib}\,S,\ \uc{Liss}\,S \qquad \text{et} \qquad \uc{Cofib}_0S,\ \uc{Cofib}\,S,\ \uc{Prop}\,S \qquad \text{sur } S.
\]
Mais l'ennui pour $p$ et $q$, c'est que les foncteurs en sens inverse que j'avais construit (et noté encore $\Phi_0$, $\Psi_0$ (mais c'étaient essentiellement $\Phi\circ r\circ C$ et $\Psi\circ r\circ c$, à ne pas confondre avec les foncteurs $\Phi_0$, $\Psi_0$ en bas de la page précédente, et marqués aussi dans le diagramme p. 8), notons les plutôt $\Phi'$, $\Psi'$, marchent bien du côté $\Cat/S$, i.e. $p\Phi'$, $q\Psi'$ reçoivent le foncteur $\mathrm{id}_{\Cat/S}$ par une $W$-flèche, mais pas du côté des catégories but - on n'arrive pas à trouver une flèche naturelle entre $\Phi'p$ resp. $\Psi'q$, et les foncteurs identiques de $\uc{Fibsc}_0S$ resp. $\uc{Cofibsc}_0S$. Je ne vois
\msp{page 53}
pas, en effet, qu'il y en ait - la difficulté, par cette voie, parait irrémédiable. Par contre, si on regarde $\Phi'$, $\Psi'$ comme foncteurs à valeurs dans $\uc{Fib}_0S$, $\uc{Cofib}_0S$ respectivement (donc dans les catégories où on n'exige pas que les flèches soient catésiennes resp. cocartésiennes), on trouve encore des $W$-morphismes % [sic] catésiennes -> cartésiennes, missing r
\[
\mathrm{id} \longrightarrow \Phi'p, \qquad \mathrm{id} \longrightarrow \psi'q, % [sic] \psi'q -> \Psi'q, lowercase psi where \Psi' is used everywhere else
\]
ce qui liquide au moins déjà le cas de toutes les catégories des diagrammes p. 8 qui sont des
\printedpage{33}
% head: A. Grothendieck : Dérivateurs, Ch. VI, Édition des Universités de Montpellier II et Paris VII
% foot: Version 9 novembre 2001
sous-catégories pleines de $\Cat/S$ (\footnote[36]{Mais c'est idiot de prendre des foncteurs si compliqués, parce qu'on veut tout ramener à $\Cat/S$. Il est de plus simple de comparer $\uc{Fibhot}\,S$ directement à $\uc{Fibsc}_0S$ et à $\uc{Cofibsc}_0S$. $\Phi_0$ et $\Psi_0$ (qui sont en fait définis sur $\uc{Fibhot}$, pas seulement sur $\uc{Parf}$), et on trouve comme ci-dessous $\mathrm{id} \xrightarrow{W} \pi_0\Phi_0$ et $\Phi_0\pi_0 \xrightarrow{W} \mathrm{id}$ (itou pour $\Psi_0$ et $\uc{Cofibsc}_0S$), où $\pi_0 : \uc{Fibsc}_0S \longrightarrow \uc{Fibhot}\,S$.}). Il ne reste à traiter que le cas des quatre catégories de structures `scindées'
\[
\uc{Fibsc}\,S,\quad \uc{Fibsc}_0S \qquad \text{et} \qquad \uc{Cofibsc}\,S,\quad \uc{Cofibsc}_0S.
\]
Bien sûr il suffit de le faire d'un côté, disons du côté des fibrés, par dualité ($S \longmapsto S^o$) cela établira le résultat symétrique. On va donc regarder d'un peu plus près le carré de foncteurs
\[
\begin{tikzcd}[column sep=huge, row sep=large]
\uc{Fibsc}\,S \arrow[d, "p"] & \uc{Fibsc}_0S \arrow[l, hook, "\alpha'"'] \arrow[d, "p_0"'] \\
\uc{Fib}\,S \arrow[u, dotted, bend left, "\varphi"] & \uc{Fib}_0S \arrow[l, hook, "\alpha"'] \arrow[u, dotted, bend right, "\varphi_0"']
\end{tikzcd}
\]
(\footnote[37]{NB Il s'agit de prouver seulement que $p$, $p_0$ induisent des équivalences pour les localisés (pour la $W$-équivalence grossière). On sait déjà qu'il en est ainsi pour $\alpha$, et ça en résultera pour $\alpha'$.}, \footnote[38]{inutile de regarder ce carré, cf. annotations marginales ci-dessus (haut de la page)}). Le foncteur
\[
\Phi : \Cat/S \longrightarrow \uc{Fibsc}\,S
\]
induit des foncteurs en sens inverse de $p$, $p_0$ sur les sous-catégories pleines $\uc{Fib}\,S$, $\uc{Fib}_0S$ de $\Cat/S$, foncteurs marqués en pointillés comme $\varphi$, $\varphi_0$.

\msp{page 54}

Il faut seulement vérifier que si $X/S$ est dans $\uc{Fib}_0S$, alors $\varphi(X) = \Phi(X)$ est dans la sous-catégorie pleine $\uc{Fibsc}_0S$ de $\uc{Fibsc}\,S$. Mais on a vu que c'est le cas, pour peu seulement que $X/S$ soit dans $\uc{Fibhot}\,S$, car les morphismes de changement de base dans $\Phi(X)$, associés aux flèches $u : s \longrightarrow t$ dans $S$
\[
u^*_{\Phi_X} : t\backslash X \longrightarrow s\backslash X
\]
sont des $W_\infty$-équivalences si $X$ est Hot-fibré sur $S$. D'autre part, pour $X$ quelconque dans $\Cat/S$, on a une flèche, fonctorielle en $X$,
\[
X \longrightarrow \Phi(X),
\]
i.e. $\mathrm{id}_{\Cat/S} \longrightarrow \pi\Phi$ où $\pi : \uc{Fibsc}\,S \longrightarrow \Cat/S$ est le foncteur canonique, qui est une $W_\infty$-équivalence (même un homotopisme), donc une $W$-équivalence. A fortiori les foncteurs $p\varphi$, $p_0\varphi_0$ donnent lieu ($\uc{Fib}\,S$ et $\uc{Fib}_0S$ étant des sous-catégories pleines de $\Cat/S$) à des flèches
\[
\mathrm{id} \longrightarrow p\varphi, \qquad \mathrm{id} \longrightarrow p_0\varphi_0.
\]

\end{document}
